% META: {"id":"1SPE-EXPO-EX-023-CDP","chapitre":"1SPE-EXPONENTIELLE","type_objet":"coup_de_pouce","exercice_id":"1SPE-EXPO-EX-023","status":"approved","extension_codes":["X1"],"programme_alignment":"OPTIONAL_EXTENSION","extension_label":"Approfondissement — Vers la Terminale","origin":"generated","status_history":["generated","audited","accepted","approved"],"release_acceptance":"RELEASE_OWNER_BATCH_ACCEPTANCE","acceptance_closure_digest":"sha256:9b3ccf9a81c5520fbb7e03b7b2d3e7bf2057a02834b6d908bf3be5f49f3b553f","acceptance_receipt_digest":"sha256:4fad86f0282e0fa4e0419cca1ad6379ae7af5a280c04a4a90880f90a1c044242"}
\begin{center}\textbf{Approfondissement — Vers la Terminale}\end{center}
\begin{coupdepouce}{1SPE-EXPO-EX-023}
\coupDePouce{1}{Rappelle les deux limites fondamentales de l'exponentielle : en $+\infty$ et en $-\infty$.}
\coupDePouce{2}{Pour la question 3, pose $X = -x$. Quand $x \to +\infty$, vers quoi tend $X$ ?}
\end{coupdepouce}
